% ===========================================================================
%  Colors, boxes and notation shared by all parts.
%  Colors match the Atlas beamer theme of the decks.
% ===========================================================================

% ulem (loaded by document-style) underlines \emph; keep italics.
\normalem

% ---- Colors (same RGB values as styles/beamerthemeAtlas.sty) ----
\definecolor{AtlasAccent}{RGB}{35, 62, 88}
\definecolor{AtlasInk}{RGB}{35, 35, 38}
\colorlet{AtlasCool}{AtlasAccent}
\definecolor{AtlasWarm}{RGB}{184, 115, 51}
\definecolor{AtlasSlate}{RGB}{108, 79, 122}
\definecolor{GoodGreen}{RGB}{46, 139, 87}
\definecolor{BadRed}{RGB}{200, 60, 50}
\definecolor{NoteGray}{RGB}{90, 90, 90}

\hypersetup{colorlinks=true, linkcolor=AtlasAccent, urlcolor=AtlasAccent, citecolor=AtlasAccent}

% ---- Boxes ----
% Every box: thin frame, light tint, small title. All boxes can break across pages.
\tcbset{
  notes/.style={enhanced, breakable, arc=2pt, boxrule=0.6pt,
    left=6pt, right=6pt, top=4pt, bottom=4pt,
    fonttitle=\bfseries\small, coltitle=white,
    before skip=7pt, after skip=7pt},
  % one centered formula (or a few), no title; gather* is set inline so the
  % box has equal padding (same trick as the decks' formulabox)
  eq/.style={notes, colback=white, colframe=AtlasCool, halign=center,
    top=6pt, bottom=6pt, before upper=\boxedgather},
}
% Tighter display skips inside every box (applied after the box font)
\tcbset{notes/.append style={before upper app={\setlength{\abovedisplayskip}{4pt}%
  \setlength{\belowdisplayskip}{4pt}\setlength{\abovedisplayshortskip}{2pt}%
  \setlength{\belowdisplayshortskip}{4pt}}}}
\newcommand{\boxedgather}{\renewenvironment{gather*}%
  {$\displaystyle\begin{gathered}}{\end{gathered}$}}
% Definition --- #1 = name
\newtcolorbox{defbox}[2][]{notes, colback=AtlasSlate!5, colframe=AtlasSlate,
  colbacktitle=AtlasSlate, title={Definition: #2}, #1}
% Theorem / lemma / corollary --- #1 = full title, e.g. "Theorem 4.1 (K&L)"
\newtcolorbox{thmbox}[2][]{notes, colback=AtlasCool!5, colframe=AtlasCool,
  colbacktitle=AtlasCool, title={#2}, #1}
% Key idea / take-away
\newtcolorbox{keybox}[1][]{notes, colback=GoodGreen!6, colframe=GoodGreen,
  colbacktitle=GoodGreen, title={Key idea}, #1}
% Limitation / pitfall
\newtcolorbox{warnbox}[1][]{notes, colback=BadRed!5, colframe=BadRed,
  colbacktitle=BadRed, title={Limitation}, #1}
% Remark / intuition: grey bar on the left, no frame
\newtcolorbox{notebox}[1][]{enhanced, breakable, frame hidden, boxrule=0pt, arc=0pt,
  colback=NoteGray!6, borderline west={2.5pt}{0pt}{NoteGray!60},
  left=8pt, right=6pt, top=3pt, bottom=3pt, before skip=6pt, after skip=6pt,
  fontupper=\small, #1}
% Warm / cool pair (value side vs. policy side; lane A vs. lane B)
\newtcolorbox{warmbox}[1][]{notes, colback=AtlasWarm!6, colframe=AtlasWarm,
  colbacktitle=AtlasWarm, #1}
\newtcolorbox{coolbox}[1][]{notes, colback=AtlasCool!6, colframe=AtlasCool,
  colbacktitle=AtlasCool, #1}
\newtcolorbox{slatebox}[1][]{notes, colback=AtlasSlate!6, colframe=AtlasSlate,
  colbacktitle=AtlasSlate, #1}
% Algorithm
\newtcolorbox{algbox}[2][]{notes, colback=white, colframe=AtlasInk,
  colbacktitle=AtlasInk, title={Algorithm: #2}, #1}
% One formula, framed
\newtcolorbox{eqbox}[1][]{eq, #1}

% Paper card at the start of a section.
% #1 title, #2 authors, #3 venue, #4 what it introduces, #5 local copy
\newcommand{\paper}[5]{%
  \begin{tcolorbox}[enhanced, arc=3pt, boxrule=0pt, colback=AtlasCool, coltext=white,
      left=10pt, right=10pt, top=6pt, bottom=6pt, before skip=6pt, after skip=8pt]
    {\large\bfseries #1\par}\vspace{2pt}
    {\small #2 \quad$\cdot$\quad #3\par}\vspace{4pt}
    {\color{white!75!AtlasCool}\rule{0.25\linewidth}{0.5pt}\par}\vspace{2pt}
    {\small\textbf{Introduces:} #4\par}
    {\scriptsize\color{white!80!AtlasCool}Local copy: \texttt{#5}\par}
  \end{tcolorbox}}

% Source line (grey, small) --- where the content comes from
\newcommand{\src}[1]{\par\noindent{\footnotesize\color{NoteGray}\textbf{Source:} #1}\par}

% Side-by-side boxes: \begin{pair} \begin{warmbox}..\end{warmbox} \begin{coolbox}..\end{coolbox} \end{pair}
\newenvironment{pair}[1][2]{\tcbset{notes/.append style={unbreakable}}%
  \begin{tcbraster}[raster columns=#1, raster equal height,
    raster column skip=8pt, raster before skip=6pt, raster after skip=6pt]}{\end{tcbraster}}

% Why needed / What changes / What it costs --- for each TRPO simplification
\newcommand{\simp}[3]{%
  \begingroup\tcbset{notes/.append style={unbreakable}}%
  \begin{tcbraster}[raster columns=3, raster equal height, raster column skip=6pt,
      raster before skip=4pt, raster after skip=8pt]
    \begin{tcolorbox}[notes, colback=AtlasAccent!4, colframe=AtlasAccent!60, colbacktitle=AtlasAccent,
        title={Why needed}, fontupper=\small]#1\end{tcolorbox}
    \begin{tcolorbox}[notes, colback=GoodGreen!4, colframe=GoodGreen!60, colbacktitle=GoodGreen,
        title={What changes}, fontupper=\small]#2\end{tcolorbox}
    \begin{tcolorbox}[notes, colback=BadRed!4, colframe=BadRed!60, colbacktitle=BadRed,
        title={What it costs}, fontupper=\small]#3\end{tcolorbox}
  \end{tcbraster}\endgroup}

% Progress strip for TRPO's five practical simplifications; #1 = current step (1-5)
\tikzset{
  ts/.style={rectangle, rounded corners=2pt, line width=0.6pt, text width=2.3cm,
    align=center, minimum height=0.55cm, inner sep=2pt, font=\scriptsize},
  ts done/.style={ts, draw=AtlasSlate!40, fill=AtlasSlate!6, text=NoteGray},
  ts cur/.style={ts, draw=AtlasCool, fill=AtlasCool!18, line width=1pt, font=\scriptsize\bfseries},
  ts todo/.style={ts, draw=AtlasSlate!25, text=AtlasSlate!45},
  % generic boxes for flow diagrams: cbox=<color>
  cbox/.style={draw=#1, fill=#1!7!white},
  fl/.style={rectangle, rounded corners=3pt, line width=0.7pt, align=center,
    inner sep=4pt, font=\small},
}
\newcommand{\trpostep}[1]{%
  \par\noindent\hfill
  \begin{tikzpicture}[baseline=(t1.base)]
    \foreach \i/\t in {1/{1.~penalty $\to$ constraint},
                       2/{2.~max $\to$ mean KL},
                       3/{3.~samples + IS},
                       4/{4.~linear / quadratic},
                       5/{5.~CG + line search}} {
      \ifnum\i<#1\relax \node[ts done] (t\i) at (\i*2.75,0) {\t};%
      \else\ifnum\i=#1\relax \node[ts cur] (t\i) at (\i*2.75,0) {\t};%
      \else \node[ts todo] (t\i) at (\i*2.75,0) {\t};\fi\fi
    }
    \foreach \i [count=\j from 2] in {1,...,4} \draw[-{Stealth[length=3pt]}, AtlasSlate!50] (t\i) -- (t\j);
  \end{tikzpicture}\hfill\null\par\vspace{2pt}}

% Ragged-right X column for tabularx
\newcolumntype{L}{>{\raggedright\arraybackslash}X}

% Check / cross marks
\newcommand{\cmark}{\leavevmode{\color{GoodGreen}\checkmark}}
\newcommand{\xmark}{\leavevmode{\color{BadRed}$\times$}}

% Part heading with a one-line summary
\newcommand{\partintro}[1]{%
  \begin{tcolorbox}[enhanced, frame hidden, colback=AtlasAccent!5, arc=2pt,
      borderline west={3pt}{0pt}{AtlasAccent}, left=8pt, top=4pt, bottom=4pt,
      before skip=4pt, after skip=10pt, fontupper=\small]#1\end{tcolorbox}}

% ---- pgfplots defaults ----
\pgfplotsset{
  notes/.style={width=7cm, height=4.6cm, axis lines=left, tick label style={font=\scriptsize},
    label style={font=\small}, legend style={font=\scriptsize, draw=none, fill=none},
    every axis plot/.append style={line width=1pt}},
}

% ---- Notation (same macros as the decks) ----
\newcommand{\E}{\mathbb{E}}
\newcommand{\PA}{\mathbb{A}}              % CPI policy advantage
\newcommand{\DKL}{D_{\mathrm{KL}}}
\newcommand{\DTV}{D_{\mathrm{TV}}}
\newcommand{\told}{\theta_{\mathrm{old}}}
\newcommand{\piold}{\pi_{\mathrm{old}}}
\newcommand{\pinew}{\pi_{\mathrm{new}}}
\newcommand{\OPT}{\mathrm{OPT}}
